\chapter{Order and chance in the integers}\label{ch:oc-integers}\label{ch:anatomy}

A deterministic set of integers can match the mean and the variance of a random model and still fail a short-interval test. This chapter fixes what ``random-like'' will mean, through the anatomy of a typical integer: normal order, friable integers, and the Dickman function. Cram\'er's model is defined in Chapter~\ref{ch:oc-models}; the distribution of $\omega(n)$ is proved in Chapter~\ref{ch:oc-divisors}. The normal order of $\omega(n)$ is the theorem of Volume~6, not a theorem proved again here.

\section{A glossary}

Three kinds of sentence appear in this volume, and they are not interchangeable.

A \emph{theorem} is proved in one of Volumes~1--14, or in this volume, and the proof is named. The normal order of $\omega(n)$, the prime number theorem, the Bombieri--Vinogradov theorem, the sine kernel, and the class-number formula are theorems. So are Buchstab's theorem, Maier's theorem, the Kubilius model, the Selberg--Delange theorem, the Bohr--Jessen law, and Selberg's central limit theorem, which this volume proves.

A \emph{model} is a probability space built to match a list of statistics. Cram\'er's coins, the Kubilius product, a Steinhaus random multiplicative function, and the Cohen--Lenstra measure are models. A model can be a theorem about its own coins and still be false about the primes. Maier's theorem is that sentence for Cram\'er's model.

A \emph{conjecture} is labelled as a conjecture in its heading. Chowla's conjecture, the Hardy--Littlewood conjectures past the second moment proved in Chapter~\ref{ch:oc-short}, the Cohen--Lenstra heuristic for class groups, and the form factor for $\alpha>1$ are conjectures. A computation that agrees with a conjecture is a computation.

The chapters own the statistics as follows. Normal order of $\omega(n)$: Volume~6. Friable integers and Dickman's delay equation: this chapter, with the permutation model cited back to Volume~11. Buchstab, Maier, and the Erd\H os--Rankin bound: Chapter~\ref{ch:oc-models}. Kubilius and Selberg--Delange: Chapter~\ref{ch:oc-divisors}. Random multiplicative functions: Chapter~\ref{ch:oc-multiplicative}. The conditional Poisson law in short intervals: Chapter~\ref{ch:oc-short}. Sieves as expectations: Chapter~\ref{ch:oc-sieves}. The Cohen--Lenstra measure: Chapter~\ref{ch:oc-classgroups}. Random isotropics: Chapter~\ref{ch:oc-ranks}. The form factor: Chapter~\ref{ch:oc-gue}. Selberg's central limit theorem: Chapter~\ref{ch:oc-values}. The $U^2$ count: Chapter~\ref{ch:oc-pseudorandom}. Chebyshev's bias, under its hypotheses: Chapter~\ref{ch:oc-bias}.

\begin{center}\itshape
Take an integer apart and look at its prime factors, its divisors, its smooth part.\\
Each organ has its law, and the laws are those of a broken stick.\\
Some of them break codes.
\end{center}

In August 1977 Martin Gardner's column in \emph{Scientific American} announced a new kind of cipher, the system of Rivest, Shamir and Adleman, and offered a challenge: a message encrypted with a public key of $129$ digits, and a prize of $100$ dollars for reading it. The column reported Ron Rivest's estimate that, with the best methods known, factoring such a number would take about $40$ quadrillion years. In April 1994 Derek Atkins, Michael Graff, Arjen Lenstra and Paul Leyland announced the factors. About six hundred volunteers had lent their computers for eight months, running the \emph{quadratic sieve}, a descendant of an algorithm that John Dixon had published in 1981. The message read: \emph{The magic words are squeamish ossifrage.}

What changed between 1977 and 1994 was not only the hardware. Dixon's algorithm and its descendants look for \emph{smooth numbers}, integers with no large prime factor, and their running time is governed by one question of probabilistic number theory: how likely is a random integer of a given size to be smooth? Volume~11, Chapter~12 answered it when the smoothness bound is a fixed power $x^{1/u}$ of the size: the proportion tends to Dickman's $\rho(u)$. Factoring needs much smaller bounds, of size $e^{\sqrt{\log x\log\log x}}$, where $u\to\infty$, and Section~\ref{sec:rankin} handles that range with a trick of Robert Rankin from 1938.

Smooth numbers are one organ of the \emph{anatomy of integers}, the study of how the prime factors of a typical integer are arranged. This chapter dissects three more: the random squares of Dixon (Section~4.2); the size of a random divisor of a random integer, which follows the arcsine law, like the time a random walk spends on one side (Section~4.3); and the number of different entries in the $N\times N$ multiplication table, which is smaller than $N^2$ by a strange power of $\log N$ (Section~4.4). Each answer is a small story about the Poisson--Dirichlet broken stick, and each uses one of the tools of Chapters~2 and~3.

Throughout, $P^+(n)$ is the largest prime factor of $n$ (with $P^+(1)=1$), $n$ is \emph{$y$-smooth} when $P^+(n)\le y$, and
\[
\Psi(x,y)=\#\{n\le x:P^+(n)\le y\},\qquad u=\frac{\log x}{\log y}.
\]

\begin{reminder}[from Volumes~2, 6 and~11]
\emph{The broken stick} (Volume~11, Chapter~12). $U_1,U_2,\dots$ are independent and uniform on $(0,1)$; $V_1=U_1$ and $V_i=U_i\prod_{j<i}(1-U_j)$ form the GEM law, and $V_{(1)}\ge V_{(2)}\ge\cdots$, their decreasing rearrangement, the Poisson--Dirichlet law. Almost surely all $V_i$ are positive and distinct and $\sum_iV_i=1$ (Lemma~12.6). The sequence $V'_i=\frac{V_{i+1}}{1-U_1}$ is GEM and independent of $U_1$.
\emph{Billingsley's theorem} (Theorem~12.17 and Exercises~12.4 and~12.5). Let $P_1(n)\ge P_2(n)\ge\cdots$ be the prime factors of $n$ with multiplicity, followed by $1$'s, and $Z_i=\frac{\log P_i(n)}{\log n}$ (with $Z_i=0$ for $n=1$). For $n$ uniform in $\{1,\dots,\lfloor x\rfloor\}$ and every $m$, $\mathbb Ef(Z_1,\dots,Z_m)\to\mathbb Ef(V_{(1)},\dots,V_{(m)})$ for every bounded Lipschitz $f$ on $[0,1]^m$.
\emph{Dickman} (Corollary~12.18). $\Psi(x,x^{1/u})\sim\rho(u)x$ for fixed $u$, and $\rho(u)\le\frac1{\lfloor u\rfloor!}$ (Proposition~12.11).
\emph{Mertens and Chebyshev} (Volume~6, Theorems~1.3 and~1.5). $\sum_{p\le x}\frac{\log p}p=\log x+O(1)$ and $\sum_{p\le x}\frac1p=\log\log x+O(1)$. The prime number theorem (Volume~6, Theorem~5.5) gives $\pi(y)\ge\frac y{2\log y}$ for $y$ large.
\emph{Units modulo $n$} (Volume~2, Theorem~3.5, Theorem~3.8 and Theorem~4.5). The Chinese remainder theorem; Hensel's lemma; $(\Z/p^e\Z)^\times$ is cyclic for odd $p$.
\end{reminder}

\section{Smooth numbers in the wide range}\label{sec:rankin}

Dickman's theorem is a statement at a fixed $u$. For the applications, $u$ grows with $x$, and $\rho(u)$ is so small that an error term $o(1)$ is useless. Rankin's idea is one line of algebra, the Chernoff bound of Volume~4 in disguise.

\begin{proposition}[Rankin's trick]\label{prop:rankin}
For $x\ge1$, $y\ge2$ and $0<\sigma\le1$,
\[
\Psi(x,y)\le x^\sigma\prod_{p\le y}\Bigl(1-\frac1{p^\sigma}\Bigr)^{-1}.
\]
\end{proposition}
\begin{proof}
For $n\le x$, $1\le(\frac xn)^\sigma$, and for $n>x$, $(\frac xn)^\sigma>0$. So
$\Psi(x,y)\le\sum_{P^+(n)\le y}(\frac xn)^\sigma=x^\sigma\sum_{P^+(n)\le y}n^{-\sigma}$, and the last sum is the finite Euler product $\prod_{p\le y}\sum_{k\ge0}p^{-k\sigma}$ (unique factorisation; every series converges since $p^{-\sigma}<1$).
\end{proof}

The bound weighs each smooth integer by $(\frac xn)^\sigma$, which is at least $1$ for the integers we count and rewards none of the others; the art is the choice of $\sigma$. With $\sigma=1$ it gives $\Psi(x,y)\ll x\log y$, which is worse than trivial; the gain comes from moving $\sigma$ below $1$, to $1-\frac{\log(u\log u)}{\log y}$.

\begin{theorem}\label{thm:rankin}
Fix $0<\sigma_0<1$. There is a constant $C=C(\sigma_0)$ such that, for $x\ge y\ge3$ with $u\ge e$ and $\log(u\log u)\le(1-\sigma_0)\log y$,
\[
\Psi(x,y)\le x\exp\bigl(-u\log u-u\log\log u+C(u+\log\log y)\bigr).
\]
\end{theorem}
\begin{proof}
Let $\eta=\frac{\log(u\log u)}{\log y}$ and $\sigma=1-\eta$, so that $\sigma_0\le\sigma<1$ and $\eta>0$ (since $u\log u\ge e$). Write $Y=\log y$.

\emph{The factor $x^\sigma$.} $x^{-\eta}=e^{-\eta uY}=e^{-u\log(u\log u)}$, so $x^\sigma=x\exp(-u\log u-u\log\log u)$.

\emph{The product.} For $0\le w\le2^{-\sigma_0}$, $-\log(1-w)=\sum_k\frac{w^k}k\le\frac w{1-w}\le c_0w$ with $c_0=(1-2^{-\sigma_0})^{-1}$. With $w=p^{-\sigma}\le2^{-\sigma_0}$,
\[
\log\prod_{p\le y}\Bigl(1-\frac1{p^\sigma}\Bigr)^{-1}\le c_0\sum_{p\le y}\frac{p^\eta}p=c_0\Bigl(\sum_{p\le y}\frac1p+\sum_{p\le y}\frac{e^{\eta\log p}-1}p\Bigr).
\]
The first sum is at most $\log\log y+C_1$ (Mertens). For the second, let $h(t)=\frac{e^{\eta t}-1}t$ for $t>0$ and $h(0)=\eta$. Since $e^{\eta t}-1$ is convex and vanishes at $0$, its slope $h$ from the origin is increasing. Let $A(t)=\sum_{p\le e^t}\frac{\log p}p$; by Volume~6, Theorem~1.3, there is $C_2\ge1$ with $|A(t)-t|\le C_2$ for all $t\ge0$ (for $t<\log2$, $A(t)=0$). By partial summation (Volume~6, Lemma~1.1, in its Stieltjes form) and $h'\ge0$,
\[
\begin{aligned}\sum_{p\le y}\frac{e^{\eta\log p}-1}p&=\int_0^Yh(t)\dd A(t)=h(Y)A(Y)-\int_0^YA(t)h'(t)\dd t\\&\le h(Y)(Y+C_2)-\int_0^Y(t-C_2)h'(t)\dd t.\end{aligned}
\]
Integrating $\int_0^Yth'(t)\dd t=Yh(Y)-\int_0^Yh$ by parts, the right side is $\int_0^Yh(t)\dd t+C_2(2h(Y)-h(0))\le\int_0^Yh(t)\dd t+2C_2h(Y)$.

Now $h(Y)=\frac{u\log u-1}{\log y}\le\frac{u\log u}{\log u}=u$, because $\log y\ge\log(u\log u)\ge\log u\ge1$. And, with $v=\eta t$ and $w=\eta Y=\log(u\log u)$,
\[
\int_0^Yh(t)\dd t=\int_0^w\frac{e^v-1}v\dd v.
\]
The integrand increases with $v$ (the same convexity argument). If $w\le2$, the integral is at most $2\cdot\frac{e^2-1}2<7$. If $w>2$, split at $1$ and at $\frac w2$: on $[0,1]$ the integrand is at most $e-1<2$; on $[1,\frac w2]$ it is at most $e^v$; on $[\frac w2,w]$ it is at most $\frac2we^v$. So the integral is at most $2+e^{w/2}+\frac2we^w$. Here $e^w=u\log u$, so $e^{w/2}=\sqrt{u\log u}\le u$ and $\frac{e^w}w=\frac{u\log u}{\log(u\log u)}\le u$. In both cases $\int_0^Yh\le7+3u$.

Collecting, the logarithm of the product is at most $c_0(\log\log y+C_1+7+(3+2C_2)u)$, and Proposition~\ref{prop:rankin} gives the theorem with $C=c_0(C_1+10+2C_2)$, since $u\ge1$.
\end{proof}

For $y=x^{1/u}$ with $u$ fixed and large, the theorem gives $\Psi(x,y)\le xe^{-u\log u(1+o(1))}$, the size $u^{-u}$ of $\rho(u)$ that de Bruijn proved in 1951; for this the hypothesis $\log(u\log u)\le(1-\sigma_0)\log y$ holds as soon as $x$ is large. But the theorem also covers $u$ as large as $\frac{y^{1-\sigma_0}}{\log y}$, far beyond any fixed power. A lower bound of the same shape is easy, by counting products of distinct small primes.

\begin{proposition}\label{prop:smoothlower}
For $x\ge y\ge2$, $\Psi(x,y)\ge\binom{\pi(y)}{\lfloor u\rfloor}$.
\end{proposition}
\begin{proof}
Let $k=\lfloor u\rfloor$. The products of $k$ distinct primes $\le y$ are $y$-smooth, at most $y^k\le y^u=x$, and pairwise distinct by unique factorisation.
\end{proof}

\begin{corollary}\label{cor:psilogpower}
For fixed $A>1$, $\Psi(x,(\log x)^A)=x^{1-1/A+o(1)}$ as $x\to\infty$.
\end{corollary}
\begin{proof}
Let $L=\log x$, $y=L^A$, so $\log y=A\log L$ and $u=\frac L{A\log L}$; then $\log u=\log L-\log\log L-\log A$, and $u\log u=\frac LA(1+o(1))$.

\emph{Upper bound.} We have $\eta=\frac{\log(u\log u)}{\log y}$ and $\log(u\log u)=\log L+O(\log\log L)$, so $\eta\to\frac1A$, and with $\sigma_0=\frac12(1-\frac1A)$ the hypothesis of Theorem~\ref{thm:rankin} holds for $x$ large, and $\log\Psi(x,y)\le L-u\log u+O(u+\log\log y)=L(1-\frac1A+o(1))$, since $u=o(L)$.

\emph{Lower bound.} Let $k=\lfloor u\rfloor$ and $h=\pi(y)\ge\frac y{2\log y}$ for $x$ large; $k\le h$. Since $\frac{h-i}{k-i}\ge\frac hk$ for $0\le i<k$, $\binom hk\ge(\frac hk)^k$, and
\[
\log\binom hk\ge k\bigl(\log y-\log\log y-\log2-\log k\bigr)\ge k\bigl((A-1)\log L-\log2\bigr),
\]
using $\log k\le\log u=\log L-\log\log L-\log A$ and $\log\log y=\log\log L+\log A$. With $k=\frac L{A\log L}(1+o(1))$ this is $L(1-\frac1A)(1+o(1))$. Apply Proposition~\ref{prop:smoothlower}.
\end{proof}

So the integers up to $x$ built from primes below $(\log x)^2$ number about $\sqrt x$, those built from primes below $(\log x)^3$ about $x^{2/3}$. Figure~\ref{fig:smooth} compares the bounds with the truth at $x=10^7$.

\begin{figure}[ht]
\centering\fig{ch04-smooth}
\caption{$\log_{10}\Psi(10^7,y)$ (charcoal dots), Rankin's bound $x^\sigma\prod_{p\le y}(1-p^{-\sigma})^{-1}$ with the best $\sigma$ (rust), the lower bound $\binom{\pi(y)}{\lfloor u\rfloor}$ of Proposition~\ref{prop:smoothlower} (ochre) and Dickman's approximation $10^7\rho(u)$ (dashed, for $u\le8.5$). Rankin's bound is too large by a factor between $3$ and $25$: small on this scale, and the price of a one-line proof. Dickman's $\rho$ is excellent for $y\ge100$ and collapses for small $y$, where $u$ is large: $10^7\rho(7)\approx9$ predicts nine $10$-smooth numbers below $10^7$, while there are $2155$. The lower bound jumps when $\lfloor u\rfloor$ drops.}
\label{fig:smooth}
\end{figure}

\begin{remark}
The truth is known with great precision. Canfield, Erd\H{o}s and Pomerance proved in 1983 that $\Psi(x,y)=xu^{-u(1+o(1))}$ uniformly for $y\ge(\log x)^{1+\varepsilon}$ and $u\to\infty$, and Hildebrand proved in 1986 that $\Psi(x,y)\sim x\rho(u)$ uniformly for $y\ge\exp((\log\log x)^{5/3+\varepsilon})$. We only need upper and lower bounds of the right exponential shape, and Rankin's trick and Proposition~\ref{prop:smoothlower} provide them; the remark is quoted for culture (see the Sources).
\end{remark}

\section{Factoring with random squares}\label{sec:dixon}

If $X^2\equiv Z^2\pmod n$ and $X\not\equiv\pm Z$, then $n$ divides $(X-Z)(X+Z)$ but neither factor, so $\gcd(X-Z,n)$ is a proper divisor of $n$ larger than $1$. Fermat and, around 1920, Maurice Kraitchik used this to factor by hand. Dixon's idea was to manufacture such congruences at random: square random residues, keep the squares that happen to be smooth, and combine them by linear algebra modulo $2$.

\emph{Dixon's algorithm.} The input is an odd integer $n\ge3$ that is not a prime power, and a bound $y\ge2$. Let $q_1<\dots<q_h$ be the primes up to $y$, $h=\pi(y)$.
\begin{enumerate}
\item Divide $n$ by $q_1,\dots,q_h$. If one divides $n$, output it and stop.
\item Repeat until $h+1$ \emph{relations} are recorded: draw $x$ uniform in $\{1,\dots,n-1\}$; if $g=\gcd(x,n)>1$, output $g$ and stop; otherwise let $s$ be the least positive residue of $x^2$ modulo $n$, and if $s$ is $y$-smooth (found by dividing by the $q_j$), record the relation $(x,s)$ with the exponent vector $v(s)\in\mathbb F_2^h$ of $s$ modulo $2$.
\item By Gaussian elimination, find a nonempty set $I$ of relations with $\sum_{i\in I}v(s_i)=0$; it exists, since $h+1$ vectors in $\mathbb F_2^h$ are dependent. Then $\prod_{i\in I}s_i=\prod_jq_j^{2f_j}$ for some integers $f_j$. Let $X=\prod_{i\in I}x_i$ and $Z=\prod_jq_j^{f_j}$, reduced modulo $n$, and output $\gcd(X-Z,n)$.
\end{enumerate}
The choice of $I$ in step~3 depends only on $s_1,\dots,s_{h+1}$. The algorithm \emph{succeeds} if its output is a proper divisor of $n$ larger than $1$; outputs of steps 1 and~2 always are.

The analysis has two parts: the probability that the final gcd is useful, which is a statement about the group of units, and the probability that a random square is smooth, which is a statement about smooth numbers.

\begin{lemma}[Dixon]\label{lem:dixon}
Let $n=p_1^{e_1}\cdots p_r^{e_r}$ be odd with $r\ge1$ distinct prime factors, all larger than $y$, let $G=(\Z/n\Z)^\times$, and let $k\ge1$ with $y^{2k}\le n$. Then:
\begin{enumerate}[label=(\alph*)]
\item the elements $\epsilon\in G$ with $\epsilon^2=1$ form a subgroup $E$ of order $2^r$, and every square in $G$ has exactly $2^r$ square roots, forming a coset of $E$;
\item the number of $x\in G$ such that the least positive residue of $x^2$ modulo $n$ is $y$-smooth is at least $\binom{\pi(y)}k^2/4^k$.
\end{enumerate}
\end{lemma}
\begin{proof}
(a) By the Chinese remainder theorem, $G\cong\prod_i(\Z/p_i^{e_i}\Z)^\times$, and each factor is cyclic of even order $p_i^{e_i-1}(p_i-1)$, so it contains exactly two solutions of $\epsilon^2=1$ (in a cyclic group of even order $2m$ with generator $g$, $g^{2a}=1$ means $m\mid a$). Hence $|E|=2^r$. If $t^2=s$, the square roots of $s$ are the $x$ with $(xt^{-1})^2=1$, that is the coset $tE$.

(b) Let $\mathcal M$ be the set of products $m=q_{j_1}\cdots q_{j_k}$ of $k$ distinct primes $\le y$, so $|\mathcal M|=\binom{\pi(y)}k$. Each $m$ is prime to $n$. To $m$ attach the sign vector $\chi(m)=\bigl((\frac m{p_1}),\dots,(\frac m{p_r})\bigr)\in\{\pm1\}^r$ of Legendre symbols, and split $\mathcal M$ into the $2^r$ classes $\mathcal M_c=\{m:\chi(m)=c\}$. If $m_1,m_2$ lie in the same class, then $(\frac{m_1m_2}{p_i})=1$ for every $i$, so $m_1m_2$ is a square modulo every $p_i$, hence modulo every $p_i^{e_i}$ (Hensel's lemma, $p_i$ odd and $p_i\nmid m_1m_2$), hence modulo $n$ (Chinese remainder theorem). Moreover $m_1m_2\le y^{2k}\le n$, and $m_1m_2\ne n$ since $n$ has no prime factor $\le y$; so $m_1m_2$ is its own least positive residue, and it is $y$-smooth.

Let $\mathcal S$ be the set of products $m_1m_2$ with $m_1,m_2$ in the same class. By Cauchy--Schwarz, the number of such ordered pairs is $\sum_c|\mathcal M_c|^2\ge\frac{|\mathcal M|^2}{2^r}$. A given $s\in\mathcal S$ has $\Omega(s)=2k$, and $(m_1,m_2)$ is determined by which $k$ of the $2k$ prime factors of $s$ (with multiplicity) go to $m_1$; so $s$ comes from at most $\binom{2k}k\le4^k$ pairs, and $|\mathcal S|\ge\frac{|\mathcal M|^2}{2^r4^k}$. By (a) each $s\in\mathcal S$ has $2^r$ square roots in $G$, and different $s$ have disjoint sets of roots. Every root $x$ of $s$ has $x^2\equiv s$ with $s$ smooth. This gives at least $2^r|\mathcal S|\ge\frac{|\mathcal M|^2}{4^k}$ elements $x$.
\end{proof}

\begin{theorem}[Dixon, 1981]\label{thm:dixon}
Let $n$ be odd and not a prime power.
\begin{enumerate}[label=(\alph*)]
\item For every $y\ge2$, Dixon's algorithm succeeds with probability at least $\frac12$.
\item With $y=\exp\bigl(\frac12\sqrt{\log n\log\log n}\bigr)$, run the algorithm again after each failure, with fresh randomness. The expected number of arithmetic operations on integers below $n^2$ until success is
\[
\exp\Bigl((4+o(1))\sqrt{\log n\log\log n}\Bigr)\qquad(n\to\infty).
\]
\end{enumerate}
\end{theorem}
\begin{proof}
(a) If the algorithm stops in step 1 or~2, it succeeds. Otherwise $n$ has $r\ge2$ prime factors, all larger than $y$, and the algorithm reaches step~3 with relations $(x_1,s_1),\dots,(x_{h+1},s_{h+1})$. Condition on $s_1,\dots,s_{h+1}$. Each draw in step~2 that is a unit is uniform in $G$; it is recorded or not according to the value of $x^2$, and by Lemma~\ref{lem:dixon}(a) the $2^r$ units with a given square are equally likely. So, conditionally, the $x_i$ are independent and $x_i$ is uniform on the coset $t_iE$ of square roots of $s_i$, where $t_i$ is a fixed root chosen as a function of $s_i$. Write $x_i=t_i\epsilon_i$, with $\epsilon_1,\dots,\epsilon_{h+1}$ independent and uniform in $E$. Then $X=T\epsilon$ with $T=\prod_{i\in I}t_i$ and $\epsilon=\prod_{i\in I}\epsilon_i$, and $\epsilon$ is uniform in $E$: for $i_0\in I$, given the other factors, $\epsilon_{i_0}$ times a fixed element is uniform. Now $I$, $T$ and $Z$ depend only on the $s_i$, and $X^2=\prod_{i\in I}s_i=Z^2$ in $G$. The algorithm fails only if $X=\pm Z$, that is $\epsilon\in\{T^{-1}Z,-T^{-1}Z\}$, which has conditional probability at most $\frac2{2^r}\le\frac12$. Otherwise, as noted at the start of the section, $\gcd(X-Z,n)$ is a proper divisor larger than $1$.

(b) Let $\ell=\log n\log\log n$, $\log y=\frac12\sqrt\ell$, $u=\frac{\log n}{\log y}=2\sqrt{\log n/\log\log n}$, $h=\pi(y)$ and $k=\lfloor\frac u2\rfloor$. Then $y^{2k}\le y^u=n$ and $y^{2k}\ge y^{u-2}=\frac n{y^2}$. In step~2, a draw either stops the algorithm or, by Lemma~\ref{lem:dixon}(b), is recorded with probability
\[
p_0\ge\frac{\binom hk^2}{4^kn}\ge\frac1n\Bigl(\frac h{2k}\Bigr)^{2k}\ge\frac1n\Bigl(\frac y{4k\log y}\Bigr)^{2k}\ge\frac{y^{2k}}n(2\log n)^{-2k}\ge y^{-2}(2\log n)^{-u},
\]
for $n$ large, using $\binom hk\ge(\frac hk)^k$, the prime number theorem $h\ge\frac y{2\log y}$, and $4k\log y\le2u\log y=2\log n$. The number of draws needed for $h+1$ relations is at most a sum of $h+1$ geometric variables of mean $\frac1{p_0}$, so its expectation is at most $\frac{h+1}{p_0}$. Each draw costs $O(h+\log n)$ operations (a gcd, a squaring, and at most $h+\log_2n$ trial divisions), step~1 costs $h$, and step~3 costs $O(h^3+h\log n)$ (Gaussian elimination on $h+1$ vectors of length $h$, then the products). Since $h\le y$, one run costs on average at most
\[
O\bigl(y\cdot y\log n\cdot y^2(2\log n)^u+y^3\log n\bigr)=O\bigl(y^4(2\log n)^{u+1}\bigr).
\]
The number of runs until success is geometric with mean at most $2$ by (a), and the runs are independent with the same cost bound. Finally $4\log y=2\sqrt\ell$ and $(u+1)\log(2\log n)=2\sqrt{\frac{\log n}{\log\log n}}(\log\log n+\log2)(1+o(1))=2\sqrt\ell(1+o(1))$, so the expected cost is $\exp((4+o(1))\sqrt\ell)$.
\end{proof}

\begin{think}
Where did the proof of part~(a) use that $n$ is not a prime power, and where that $n$ is odd? What happens to Dixon's algorithm for $n=p^2$, and for $n=2p$? (Exercise~\ref{ex:primepower}.)
\end{think}

The choice of $y$ is a balance: a larger $y$ makes smooth squares more frequent but requires more of them, and more trial divisions. With $\log y=a\sqrt\ell$ the two costs are $e^{4a\sqrt\ell}$ and $e^{\sqrt\ell/a}$, and $4a+\frac1a$ is smallest at $a=\frac12$. A running time $e^{c\sqrt{\log n\log\log n}}$ is \emph{subexponential}: smaller than any power $n^\varepsilon$, larger than any power of $\log n$, the number of digits. Dixon's algorithm was the first factoring method with a proved subexponential running time.

\begin{remark}
Refinements of the same idea (better choices of $y$, stopping trial division early, sparse linear algebra) lower the constant $4$. The quadratic sieve of Carl Pomerance (1981) replaces the random $x$ by $x$ near $\sqrt n$, so that $x^2-n$ is small and more often smooth, and finds the smooth values by a sieve instead of trial division; its heuristic running time is $e^{(1+o(1))\sqrt{\log n\log\log n}}$. The number field sieve (around 1990) reaches the heuristic $e^{(c+o(1))(\log n)^{1/3}(\log\log n)^{2/3}}$ with $c=(64/9)^{1/3}=1.92\ldots$, and holds the factoring records today. These running times are heuristic: they assume that the numbers produced behave like random numbers of their size with respect to smoothness. The best \emph{proved} bound, $e^{(1+o(1))\sqrt{\log n\log\log n}}$, is due to H.~W.~Lenstra and Pomerance (1992), with a different algorithm. See the Sources.
\end{remark}

\section{A random divisor of a random integer}\label{sec:divisors}

Pick $n$ at random up to $x$, then a divisor $d$ of $n$ at random among its $\tau(n)$ divisors. Is $d$ typically small, near $\sqrt n$, or large? The pairing $d\leftrightarrow\frac nd$ shows that $\frac{\log d}{\log n}$ is symmetric about $\frac12$, but says nothing about its shape. In 1979 Jean-Marc Deshouillers, François Dress and Gérald Tenenbaum found the answer: $\frac{\log d}{\log n}$ follows the \emph{arcsine law}, with density $\frac1{\pi\sqrt{t(1-t)}}$, the law of the fraction of time that a random walk spends on the positive side (Volume~4, Theorem~7.7; Volume~11, Theorem~3.17). A random divisor is far more likely to be near $1$ or near $n$, on the logarithmic scale, than near $\sqrt n$.

The reason is the broken stick. A uniform divisor of a squarefree $n$ keeps each prime factor independently with probability $\frac12$, so $\frac{\log d}{\log n}$ is a random sub-sum of the pieces $\frac{\log p}{\log n}$, which follow the Poisson--Dirichlet law.

\begin{lemma}\label{lem:gemarcsine}
Let $(V_i)$ have the GEM law, and let $\varepsilon_1,\varepsilon_2,\dots$ be independent Bernoulli$(\frac12)$ variables, independent of $(V_i)$. Then $S=\sum_i\varepsilon_iV_i$ has the arcsine law: for $0\le t\le1$, $P(S\le t)=\frac2\pi\arcsin\sqrt t$. The same holds for $\sum_i\varepsilon_iV_{(i)}$.
\end{lemma}
\begin{proof}
Since $V_1=U_1$ and $V_{i+1}=(1-U_1)V'_i$,
\[
S=\varepsilon_1U_1+(1-U_1)S',\qquad S'=\sum_i\varepsilon_{i+1}V'_i,
\]
where $S'$ has the law of $S$ and is independent of $(\varepsilon_1,U_1)$. Let $s_m=\mathbb ES^m$, so $s_0=1$. By the binomial formula and independence,
\[
s_m=\sum_{k=0}^m\binom mk\mathbb E[\varepsilon_1^k]\,\mathbb E[U_1^k(1-U_1)^{m-k}]\,s_{m-k}.
\]
Here $\mathbb E[\varepsilon_1^k]=\frac12$ for $k\ge1$ and $1$ for $k=0$, and $\binom mk\mathbb E[U_1^k(1-U_1)^{m-k}]=\binom mk\frac{k!(m-k)!}{(m+1)!}=\frac1{m+1}$ (the Beta integral). So $(m+1)s_m=s_m+\frac12\sum_{j<m}s_j$, that is
\[
2ms_m=\sum_{j=0}^{m-1}s_j\qquad(m\ge1).
\]
This recursion determines all $s_m$ from $s_0=1$. The moments $a_m=\binom{2m}m4^{-m}$ of the arcsine law satisfy it: $a_0=1$, and subtracting the recursion at $m-1$ from the one at $m$ shows that it is equivalent to $2ma_m=(2m-1)a_{m-1}$, which holds since $\frac{a_m}{a_{m-1}}=\frac{2m(2m-1)}{4m^2}$. That $a_m$ are the arcsine moments: with $t=\sin^2\theta$,
\[
\int_0^1\frac{t^m\dd t}{\pi\sqrt{t(1-t)}}=\frac2\pi\int_0^{\pi/2}\sin^{2m}\theta\dd\theta=\binom{2m}m4^{-m}
\]
by Wallis's formula. Two laws on $[0,1]$ with the same moments are equal, since polynomials are dense in the continuous functions on $[0,1]$ (Weierstrass; Volume~7, Corollary~4.20). The distribution function of the arcsine law is $\frac2\pi\arcsin\sqrt t$ by the substitution $t=\sin^2\theta$.

For the last statement, the decreasing rearrangement is $V_{(i)}=V_{\pi(i)}$ for a random permutation $\pi$ that depends only on $(V_i)$ (it is well defined since the $V_i$ are distinct). Given $(V_i)$, the sequence $(\varepsilon_{\pi^{-1}(i)})$ is still independent Bernoulli$(\frac12)$, so $\sum_i\varepsilon_iV_{(i)}=\sum_i\varepsilon_{\pi^{-1}(i)}V_i$ has the law of $S$.
\end{proof}

\begin{theorem}[Deshouillers, Dress and Tenenbaum, 1979]\label{thm:ddt}
Let $n$ be uniform in $\{1,\dots,\lfloor x\rfloor\}$ and, given $n$, let $d$ be uniform among the divisors of $n$; put $D_x=\frac{\log d}{\log n}$ (and $D_x=0$ if $n=1$). Then for every $t\in[0,1]$,
\[
P(D_x\le t)\longrightarrow\frac2\pi\arcsin\sqrt t\qquad(x\to\infty).
\]
\end{theorem}
\begin{proof}
\emph{The divisor, prime by prime.} If $n=\prod_pp^{a_p}$, a uniform divisor is $\prod_pp^{b_p}$ with $b_p$ independent and uniform in $\{0,\dots,a_p\}$, since the divisors correspond to the exponent vectors $(b_p)$. Fix $m\ge1$, let $Z_i=\frac{\log P_i(n)}{\log n}$ as in the reminder, and let $A_m$ be the event $Z_1>Z_2>\dots>Z_{m+1}$. On $A_m$ the primes $P_1(n),\dots,P_m(n)$ are distinct and each divides $n$ exactly once, so, given $n$, their exponents $\varepsilon_i=b_{P_i(n)}$ are independent Bernoulli$(\frac12)$ variables, independent of the other exponents, and
\[
D_x=\sum_{i\le m}\varepsilon_iZ_i+R,\qquad0\le R\le1-\sum_{i\le m}Z_i,
\]
since the rest of $d$ divides $n/(P_1(n)\cdots P_m(n))$, whose logarithm is $\log n\,(1-\sum_{i\le m}Z_i)$. (Recall that $\sum_iZ_i=1$ for $n\ge2$.) Off $A_m$, let $\varepsilon_1,\dots,\varepsilon_m$ be fresh independent Bernoulli$(\frac12)$ variables. Then, given $n$, the $\varepsilon_i$ are always independent Bernoulli$(\frac12)$, so they are independent of $n$. Put $T=\sum_{i\le m}\varepsilon_iZ_i$.

\emph{Comparison.} Let $f:\R\to[-1,1]$ be $1$-Lipschitz. Then
\[
\bigl|\mathbb Ef(D_x)-\mathbb Ef(T)\bigr|\le2P(A_m^c)+\mathbb E\Bigl[1-\sum_{i\le m}Z_i\Bigr].
\]
The function $g(z_1,\dots,z_m)=\mathbb E_\varepsilon f(\sum\varepsilon_iz_i)$ is bounded and Lipschitz, so by Billingsley's theorem $\mathbb Ef(T)=\mathbb Eg(Z_1,\dots,Z_m)\to\mathbb Eg(V_{(1)},\dots,V_{(m)})$, and likewise $\mathbb E[1-\sum_{i\le m}Z_i]\to r_m=\mathbb E[1-\sum_{i\le m}V_{(i)}]$. With $S_\infty=\sum_{i\ge1}\varepsilon_iV_{(i)}$, extending the $\varepsilon_i$ to all $i$, $|\mathbb Eg(V_{(1)},\dots,V_{(m)})-\mathbb Ef(S_\infty)|\le r_m$.

\emph{The event $A_m$.} For $\eta>0$, let $\phi_\eta(z_1,\dots,z_{m+1})=\min\bigl(1,\sum_{i\le m}(1-\frac{z_i-z_{i+1}}\eta)^+\bigr)$: it is bounded, Lipschitz, equal to $1$ off $A_m$ (some $z_i-z_{i+1}\le0$) and at most $\sum_{i\le m}\mathbf 1\{z_i-z_{i+1}<\eta\}$. By Billingsley's theorem with $m+1$ coordinates,
\[
\limsup_xP(A_m^c)\le\lim_x\mathbb E\phi_\eta(Z_1,\dots,Z_{m+1})=\mathbb E\phi_\eta(V_{(1)},\dots,V_{(m+1)})\le\sum_{i\le m}P\bigl(V_{(i)}-V_{(i+1)}<\eta\bigr).
\]
As $\eta\to0$ the right side tends to $\sum_iP(V_{(i)}=V_{(i+1)})=0$, the $V_i$ being almost surely distinct (Volume~11, Lemma~12.6). So $P(A_m^c)\to0$.

\emph{Conclusion.} Collecting, $\limsup_x|\mathbb Ef(D_x)-\mathbb Ef(S_\infty)|\le2r_m$ for every $m$. Since $\sum_iV_{(i)}=1$ almost surely, $r_m\to0$ by dominated convergence, so $\mathbb Ef(D_x)\to\mathbb Ef(S_\infty)$ for every bounded Lipschitz $f$ (scale $f$ to reduce to the case above). By Lemma~\ref{lem:gemarcsine}, $S_\infty$ has the arcsine distribution function $F(t)=\frac2\pi\arcsin\sqrt t$, which is continuous. For $t\in[0,1]$ and $\delta>0$, let $f$ be $1$ on $(-\infty,t]$, $0$ on $[t+\delta,\infty)$ and linear between; then $\limsup P(D_x\le t)\le\lim\mathbb Ef(D_x)\le F(t+\delta)$. Similarly, with $f$ equal to $1$ on $(-\infty,t-\delta]$ and $0$ on $[t,\infty)$, $\liminf P(D_x\le t)\ge F(t-\delta)$. Let $\delta\to0$.
\end{proof}

\begin{figure}[ht]
\centering\fig{ch04-arcsine}
\caption{The distribution function of $\frac{\log d}{\log n}$ for $n\le10^6$ uniform and $d$ a uniform divisor of $n$ (charcoal), against the arcsine law $\frac2\pi\arcsin\sqrt t$ (dashed). The largest gap, $0.11$, is at $t=0$: the atom $P(d=1)=\mathbb E\frac1{\tau(n)}=0.151$, which tends to $0$ only like $(\log x)^{-1/2}$ (Exercise~\ref{ex:invtau}).}
\label{fig:arcsine}
\end{figure}

\begin{beauty}
The same law governs the fraction of time a fair coin-tossing game spends in the lead (Volume~4, Theorem~7.7) and the size of a random divisor. Both are random sub-sums of a broken stick with fair signs. For the divisor, the pieces are the prime factors, which follow the Poisson--Dirichlet law, and each is kept or dropped with probability $\frac12$. For the game, in the Brownian limit, the pieces are the excursions away from $0$, each spent on one side or the other with probability $\frac12$; their lengths follow a different broken-stick law, the Poisson--Dirichlet law with parameter $\frac12$ of Pitman and Yor. Two different sticks, one arcsine law: Lemma~\ref{lem:gemarcsine} explains the divisors in one line of moments, and the coincidence with the walk is a small miracle of the two families.
\end{beauty}

So with probability $\frac2\pi\arcsin\sqrt{0.1}=0.205$ a random divisor of a random integer is below $n^{1/10}$, and with the same probability above $n^{9/10}$, while it lies between $n^{0.45}$ and $n^{0.55}$ with probability only $0.064$.

\section{The multiplication table}\label{sec:table}

How many different numbers appear in the $N\times N$ multiplication table? Let
\[
D(N)=\#\{ab:1\le a,b\le N\}.
\]
For $N=10$ the table has $100$ entries but only $D(10)=42$ different values. Paul Erd\H{o}s showed in 1955 that $D(N)=o(N^2)$, and in 1960 that $D(N)=N^2(\log N)^{-\delta+o(1)}$ with
\[
\delta=1-\frac{1+\log\log2}{\log2}=0.08607\ldots
\]
The heuristic is a mismatch of two Erd\H{o}s--Kac laws. A typical $n\le N^2$ has about $\log\log N$ prime factors. A typical product $ab$ with $a,b\le N$ has about $2\log\log N$ of them, twice as many: so most products are atypical integers, and an atypical integer is the product of many pairs. Counting integers by their number of prime factors, which is the business of Chapter~\ref{ch:oc-divisors}, makes this precise. The upper bound needs Corollary~\ref{cor:bigomegasd} at two values of $z$.

\begin{theorem}[Erd\H{o}s, 1960]\label{thm:multtable}
For $N\ge3$, $D(N)\ll\dfrac{N^2}{(\log N)^\delta}$.
\end{theorem}
\begin{proof}
Let $\lambda=\log\log N$, $\kappa=\frac1{\log2}$, $k_0=\lceil\kappa\lambda\rceil$, and recall $\Omega(ab)=\Omega(a)+\Omega(b)$. Every product $m=ab$ either has $\Omega(m)\ge k_0$ or comes from a pair with $\Omega(a)+\Omega(b)<k_0$. So
\[
D(N)\le\#\{m\le N^2:\Omega(m)\ge k_0\}+\#\{(a,b)\in[1,N]^2:\Omega(a)+\Omega(b)<k_0\}.
\]
By Corollary~\ref{cor:bigomegasd}, for fixed real $0<z\le\frac32$ and $X\ge3$, $\sum_{n\le X}z^{\Omega(n)}\le C_zX(\log X)^{z-1}$.

\emph{First term.} With $z_1=\kappa=1.4427\ldots\le\frac32$, $\mathbf 1\{\Omega(m)\ge k_0\}\le z_1^{\Omega(m)-k_0}$, so the first term is at most $z_1^{-k_0}C_{z_1}N^2(2\log N)^{z_1-1}\ll N^2\exp\bigl(\lambda(z_1-1-\kappa\log z_1)\bigr)$, since $z_1^{-k_0}\le z_1^{-\kappa\lambda}$.

\emph{Second term.} With $z_2=\frac\kappa2=0.7213\ldots<1$, $\mathbf 1\{\Omega(a)+\Omega(b)<k_0\}\le z_2^{\Omega(a)+\Omega(b)-k_0}$, so the second term is at most $z_2^{-k_0}\bigl(\sum_{a\le N}z_2^{\Omega(a)}\bigr)^2\ll z_2^{-\kappa\lambda-1}N^2(\log N)^{2(z_2-1)}\ll N^2\exp\bigl(\lambda(2z_2-2-\kappa\log z_2)\bigr)$.

\emph{The exponents.} The first is $\kappa-1-\kappa\log\kappa$. The second is $\kappa-2-\kappa\log\kappa+\kappa\log2=\kappa-1-\kappa\log\kappa$, since $\kappa\log2=1$: the two are equal, and this is why $\kappa$ was chosen. Finally $\kappa-1-\kappa\log\kappa=\frac{1+\log\log2}{\log2}-1=-\delta$, because $\log\kappa=-\log\log2$. So both terms are $\ll N^2e^{-\delta\lambda}=N^2(\log N)^{-\delta}$.
\end{proof}

Where does the choice of $z_1$ and $z_2$ come from? For a threshold $k_0=\alpha\lambda$, the best $z$ in each Chernoff bound is $z_1=\alpha$ and $z_2=\frac\alpha2$, which give the exponents $-Q^*(\alpha)$ and $-2Q^*(\frac\alpha2)$ of Chapter~\ref{ch:erdoskac}; and $\alpha$ is chosen to make them equal. The integers $m$ with $\Omega(m)\approx\kappa\log\log N$ are the bottleneck: by the same Chernoff bounds they form a proportion about $(\log N)^{-Q^*(\kappa)}$ of the integers up to $N^2$, and the pairs $(a,b)$ that produce them about the same proportion of all pairs; and $Q^*(\kappa)=\kappa\log\kappa-\kappa+1=\delta$.

\begin{figure}[ht]
\centering\fig{ch04-table}
\caption{$D(N)/N^2$ for $N=2,4,\dots,2^{13}$ (charcoal), computed by marking the products in a boolean array, against Ford's order of magnitude $c(\log N)^{-\delta}(\log\log N)^{-3/2}$ (dashed, with $c=0.89$ fitted at $N=2^{13}$). The decay is real but extremely slow: $D(N)/N^2$ is still $0.227$ at $N=8192$, and the dashed curve reaches $0.1$ only near $N=10^{15}$, a table with $10^{30}$ entries.}
\label{fig:table}
\end{figure}

\begin{remark}
Erd\H{o}s's lower bound $D(N)\ge N^2(\log N)^{-\delta-o(1)}$ comes from counting products $ab$ with $\Omega(a)$ and $\Omega(b)$ near $\frac\kappa2\lambda$ and showing that they are rarely repeated. The exact order of magnitude, $D(N)\asymp\frac{N^2}{(\log N)^\delta(\log\log N)^{3/2}}$, was found by Kevin Ford in 2008, after work of Tenenbaum, as part of a study of the integers with a divisor in a given interval. Figure~\ref{fig:table} shows that the factor $(\log\log N)^{-3/2}$ matters: without it the curve $(\log N)^{-\delta}$ is still above $0.8$ at $N=2^{13}$.
\end{remark}

\section{Exercises}

\begin{exercise}[$\star$]\label{ex:psi23}
Show that $\Psi(x,2)=\lfloor\frac{\log x}{\log2}\rfloor+1$ for $x\ge1$, and that $\Psi(x,3)=\frac{(\log x)^2}{2\log2\log3}+O(\log x)$.
\end{exercise}

\begin{exercise}[$\star\star$]\label{ex:psilog}
Show that $\Psi(x,\log x)=x^{o(1)}$. (Count exponent vectors.) Compare with Corollary~\ref{cor:psilogpower}.
\end{exercise}

\begin{exercise}[$\star$]\label{ex:tablelower}
Show that $D(N)\gg\frac{N^2}{\log N}$, by looking at the products $ab$ with $a\le\frac N2$ and $b$ a prime in $(\frac N2,N]$.
\end{exercise}

\begin{exercise}[$\star\star$]\label{ex:primepower}
(a) Why does Dixon's method fail when $n=p^e$ is a prime power? (b) Give an algorithm that decides whether $n$ is a perfect power $m^k$ with $k\ge2$, and finds $m$ and $k$, in a number of arithmetic operations polynomial in $\log n$.
\end{exercise}

\begin{exercise}[$\star\star$]\label{ex:invtau}
Show, with Theorem~\ref{thm:sd}, that $\sum_{n\le x}\frac1{\tau(n)}=\frac{Cx}{\sqrt{\log x}}\bigl(1+O(\frac1{\log x})\bigr)$, where
\[
C=\frac1{\sqrt\pi}\prod_p\sqrt{1-\frac1p}\;p\log\frac p{p-1}=0.5468\ldots
\]
Deduce that, in Theorem~\ref{thm:ddt}, $P(d=1)\to0$ only like $(\log x)^{-1/2}$.
\end{exercise}

\begin{exercise}[$\star\star\star$]\label{ex:betalaw}
Let $0<\theta<1$, and in Lemma~\ref{lem:gemarcsine} let the $\varepsilon_i$ be Bernoulli$(\theta)$. Show that $S$ has the Beta$(\theta,1-\theta)$ law, of density $\frac{\sin\pi\theta}\pi t^{\theta-1}(1-t)^{-\theta}$ on $(0,1)$. Interpret it for divisors chosen by keeping each prime factor of a squarefree integer with probability $\theta$.
\end{exercise}

\section{Hints}

\begin{hint}{ex:psi23}
The $3$-smooth $n\le x$ are the $2^a3^b$ with $a\log2+b\log3\le\log x$: lattice points in a triangle.
\end{hint}
\begin{hint}{ex:psilog}
A $y$-smooth $n\le x$ is determined by its exponents $(e_p)_{p\le y}$ with $\sum e_p\le\frac{\log x}{\log2}$; count with stars and bars, and use $\pi(y)\ll\frac y{\log y}$.
\end{hint}
\begin{hint}{ex:tablelower}
If $ab=a'b'$ with $b,b'$ primes larger than $\frac N2\ge a,a'$, then $b=b'$.
\end{hint}
\begin{hint}{ex:primepower}
(a) Count the square roots of $1$. (b) For each $k\le\log_2n$, find $\lfloor n^{1/k}\rfloor$ by bisection.
\end{hint}
\begin{hint}{ex:invtau}
$\frac1\tau$ is multiplicative with value $\frac1{k+1}$ at $p^k$; its Dirichlet series is $\zeta(s)^{1/2}G(s)$ where the local factors of $G$ are $1+O(p^{-2\sigma})$.
\end{hint}
\begin{hint}{ex:betalaw}
The same recursion becomes $ms_m=\theta\sum_{j<m}s_j$.
\end{hint}

\section{Solutions}

\begin{solution}{ex:psi23}
\idea Smooth numbers for a fixed set of primes are lattice points under a hyperplane, and lattice points in a nice region are counted by its area.
\writeup The $2$-smooth $n\le x$ are $2^a$ with $0\le a\le\frac{\log x}{\log2}$: there are $\lfloor\frac{\log x}{\log2}\rfloor+1$. The $3$-smooth ones are $2^a3^b$ with $(a,b)\in\N^2$ in the triangle $\Delta=\{a,b\ge0,\ a\log2+b\log3\le L\}$, $L=\log x$. For each $b$ with $0\le b\le\frac L{\log3}$, there are $\lfloor\frac{L-b\log3}{\log2}\rfloor+1$ values of $a$, which is $\frac{L-b\log3}{\log2}+O(1)$. Summing over the $O(L)$ values of $b$ and comparing the sum with the integral $\int_0^{L/\log3}\frac{L-t\log3}{\log2}\dd t=\frac{L^2}{2\log2\log3}$ (a decreasing function, so sum and integral differ by at most its maximum $\frac L{\log2}$) gives $\Psi(x,3)=\frac{L^2}{2\log2\log3}+O(L)$. In general, $\Psi(x,y)$ for fixed $y$ is $\frac{L^{\pi(y)}}{\pi(y)!\prod_{p\le y}\log p}(1+o(1))$, the volume of a simplex.
\end{solution}

\begin{solution}{ex:psilog}
\idea Smooth numbers are few when there are few primes: count the exponent vectors.
\writeup Let $L=\log x$, $y=L$ and $k=\pi(y)$. A $y$-smooth $n\le x$ is $\prod_{p\le y}p^{e_p}$ with $\sum e_p\le\sum e_p\log_2p\le\frac L{\log2}$. The number of $k$-tuples of natural numbers with sum at most $M=\lfloor\frac L{\log2}\rfloor$ is $\binom{M+k}k\le(\frac{e(M+k)}k)^k$. Now $k\le\frac{CL}{\log L}$ (Chebyshev), and $t\mapsto t\log\frac{e(M+t)}t$ is increasing, so
\[
\log\Psi(x,L)\le\frac{CL}{\log L}\log\Bigl(e+\frac{e\log L}{C\log2}\Bigr)=O\Bigl(\frac{L\log\log L}{\log L}\Bigr)=o(L).
\]
So $\Psi(x,\log x)=x^{o(1)}$, in agreement with the limit $1-\frac1A\to0$ of Corollary~\ref{cor:psilogpower} as $A\to1$. (In fact $\log\Psi(x,\log x)\sim\log4\cdot\frac{\log x}{\log\log x}$, by a finer count.)
\end{solution}

\begin{solution}{ex:tablelower}
\idea Products with one large prime factor are unique: the prime can be read off.
\writeup Let $a,a'\le\frac N2$ and $b,b'$ primes in $(\frac N2,N]$ with $ab=a'b'$. The prime $b$ divides $a'b'$ but not $a'$ (as $1\le a'<b$), so $b\mid b'$ and $b=b'$, then $a=a'$. So these products are distinct, and $D(N)\ge\lfloor\frac N2\rfloor(\pi(N)-\pi(\frac N2))$. By the prime number theorem, $\pi(N)-\pi(\frac N2)\sim\frac N{2\log N}$, so $D(N)\ge\frac{N^2}{5\log N}$ for $N$ large.
\end{solution}

\begin{solution}{ex:primepower}
\idea Dixon's method only ever learns a square root of $1$; when $n$ is a prime power the only ones are $\pm1$.
\writeup (a) If $n=p^e$ with $p$ odd, $G$ is cyclic (Volume~2, Theorem~4.5) and $E=\{\pm1\}$, so $r=1$ in Lemma~\ref{lem:dixon}. In step~3, $X^2=Z^2$ gives $(XZ^{-1})^2=1$, hence $X=\pm Z$ always, and the gcd is $1$ or $n$: the method never succeeds. (b) If $n=m^k$ with $m\ge2$ then $k\le\log_2n$. For each $k$ from $2$ to $\lfloor\log_2n\rfloor$, find the largest integer $m$ with $m^k\le n$ by bisection on $[1,n]$: $\log_2n+1$ steps, each computing $m^k$ by repeated squaring with $O(\log k)$ multiplications (stop as soon as a partial product exceeds $n$, so all numbers stay below $n^2$). Then test $m^k=n$. The total is $O((\log n)^2\log\log n)$ operations. A prime power is detected in this way, and $m$ can be tested in turn; an odd $n$ that is not a perfect power is either prime (detected by a primality test) or has $r\ge2$ prime factors, and Dixon's algorithm applies.
\end{solution}

\begin{solution}{ex:invtau}
\idea $\frac1{\tau(n)}$ is $(\frac12)^{\omega(n)}$ for squarefree $n$, so it behaves like $\zeta^{1/2}$.
\writeup $f=\frac1\tau$ is multiplicative with $f(p^k)=\frac1{k+1}$, and $0<f\le1$. For $\sigma>1$, $F(s)=\sum f(n)n^{-s}=\prod_p\phi(p^{-s})$ with $\phi(w)=\sum_{k\ge0}\frac{w^k}{k+1}=\frac{-\log(1-w)}w$. Let $G=F\zeta^{-1/2}$, with local factors $g_p(s)=\psi(p^{-s})$, $\psi(w)=\phi(w)(1-w)^{1/2}$. For $|w|<1$, $\psi(w)=(1+\frac w2+\frac{w^2}3+\cdots)(1-\frac w2-\frac{w^2}8-\cdots)=1-\frac{w^2}{24}+O(|w|^3)$; more precisely $\psi$ is holomorphic on $|w|<1$ with $\psi(0)=1$, $\psi'(0)=0$, so $|\psi(w)-1|\le C|w|^2$ for $|w|\le2^{-3/4}$. Hence the product $\prod_pg_p(s)$ converges absolutely and uniformly on $\sigma\ge\frac34$, $G$ is holomorphic and bounded there, exactly as in Lemma~\ref{lem:eulerprod}, and $U\subset\{\sigma>\frac34\}$. Theorem~\ref{thm:sd} with $z=\frac12$, $q=2$, $K=1$ gives $\sum_{n\le x}\frac1{\tau(n)}=x(\log x)^{-1/2}(\frac{G(1)}{\Gamma(1/2)}+O(\frac1{\log x}))$, and $G(1)=\prod_p(1-\frac1p)^{1/2}\phi(\frac1p)=\prod_p(1-\frac1p)^{1/2}p\log\frac p{p-1}$. Numerically the product is $0.9692\ldots$ and $C=0.5468\ldots$, a constant already computed by Ramanujan (a direct computation gives it). In Theorem~\ref{thm:ddt}, $P(d=1)=\mathbb E\frac1{\tau(n)}=\frac1{\lfloor x\rfloor}\sum_{n\le x}\frac1{\tau(n)}\sim\frac C{\sqrt{\log x}}$: at $x=10^6$ this predicts $0.147$, while the true value is $0.151$.
\end{solution}

\begin{solution}{ex:betalaw}
\idea The broken-stick recursion only sees the mean $\theta$ of $\varepsilon_1$, and the Beta law is the one whose moments obey the resulting recursion.
\writeup As in the proof of Lemma~\ref{lem:gemarcsine}, now with $\mathbb E\varepsilon_1^k=\theta$ for $k\ge1$, $(m+1)s_m=s_m+\theta\sum_{j<m}s_j$, that is $ms_m=\theta\sum_{j<m}s_j$, and by differencing $ms_m=(m-1+\theta)s_{m-1}$, so $s_m=\frac{\theta(\theta+1)\cdots(\theta+m-1)}{m!}$. The Beta$(\theta,1-\theta)$ law has moments $\frac{B(\theta+m,1-\theta)}{B(\theta,1-\theta)}=\frac{\Gamma(\theta+m)\Gamma(1)}{\Gamma(\theta)\Gamma(m+1)}$, the same numbers, and its normalising constant is $B(\theta,1-\theta)=\Gamma(\theta)\Gamma(1-\theta)=\frac\pi{\sin\pi\theta}$ (Euler's reflection formula, Volume~5). Moments determine laws on $[0,1]$, as in the lemma. For $\theta=\frac12$ we recover the arcsine law. The proof of Theorem~\ref{thm:ddt} applies verbatim: if each prime factor of $n$ that divides it exactly once is kept with probability $\theta$ (and the others are treated in any way), $\frac{\log d}{\log n}$ tends to Beta$(\theta,1-\theta)$. As $\theta\to0$ the law piles up near $0$, as it should.
\end{solution}

\section*{Sources and further reading}

\begin{itemize}
\item \emph{The challenge of 1977.} M.~Gardner, ``Mathematical games: a new kind of cipher that would take millions of years to break'', \emph{Scientific American} 237 (August 1977), 120--124, for the challenge and Rivest's estimate. D.~Atkins, M.~Graff, A.~K.~Lenstra and P.~C.~Leyland, ``The magic words are squeamish ossifrage'', in \emph{Advances in Cryptology, ASIACRYPT '94}, Lecture Notes in Computer Science 917, Springer, 1995, 263--277, for the factorisation, the volunteers and the eight months.
\item \emph{Rankin's trick.} R.~A.~Rankin, ``The difference between consecutive prime numbers'', \emph{J.\ London Math.\ Soc.}\ 13 (1938), 242--247. N.~G.~de Bruijn, ``The asymptotic behaviour of a function occurring in the theory of primes'', \emph{J.\ Indian Math.\ Soc.}\ 15 (1951), 25--32.
\item \emph{Smooth numbers.} E.~R.~Canfield, P.~Erd\H{o}s and C.~Pomerance, ``On a problem of Oppenheim concerning `factorisatio numerorum''', \emph{J.\ Number Theory} 17 (1983), 1--28. A.~Hildebrand, ``On the number of positive integers $\le x$ and free of prime factors $>y$'', \emph{J.\ Number Theory} 22 (1986), 289--307. The surveys: A.~Hildebrand and G.~Tenenbaum, ``Integers without large prime factors'', \emph{J.\ Théor.\ Nombres Bordeaux} 5 (1993), 411--484, and A.~Granville, ``Smooth numbers: computational number theory and beyond'', in \emph{Algorithmic Number Theory}, MSRI Publications 44, Cambridge University Press, 2008, 267--323.
\item \emph{Factoring.} J.~D.~Dixon, ``Asymptotically fast factorization of integers'', \emph{Math.\ Comp.}\ 36 (1981), 255--260. C.~Pomerance, ``A tale of two sieves'', \emph{Notices Amer.\ Math.\ Soc.}\ 43 (1996), 1473--1485, for the history of the quadratic and number field sieves, and Kraitchik's role. H.~W.~Lenstra and C.~Pomerance, ``A rigorous time bound for factoring integers'', \emph{J.\ Amer.\ Math.\ Soc.}\ 5 (1992), 483--516. R.~Crandall and C.~Pomerance, \emph{Prime Numbers: A Computational Perspective}, 2nd ed., Springer, 2005, Chapter~6.
\item \emph{Divisors.} J.-M.~Deshouillers, F.~Dress and G.~Tenenbaum, ``Lois de répartition des diviseurs, 1'', \emph{Acta Arith.}\ 34 (1979), 273--285. Our proof through the broken stick follows the probabilistic point of view of P.~Billingsley, ``On the distribution of large prime divisors'', \emph{Period.\ Math.\ Hungar.}\ 2 (1972), 283--289.
\item \emph{Excursions as a broken stick.} J.~Pitman and M.~Yor, ``The two-parameter Poisson--Dirichlet distribution derived from a stable subordinator'', \emph{Ann.\ Probab.}\ 25 (1997), 855--900.
\item \emph{The multiplication table.} P.~Erd\H{o}s, ``Some remarks on number theory'' (in Hebrew), \emph{Riveon Lematematika} 9 (1955), 45--48, and ``An asymptotic inequality in the theory of numbers'' (in Russian), \emph{Vestnik Leningrad.\ Univ.}\ 15 (1960), 41--49. K.~Ford, ``The distribution of integers with a divisor in a given interval'', \emph{Ann.\ of Math.}\ 168 (2008), 367--433.
\item \emph{The anatomy of integers.} J.-M.~De Koninck, A.~Granville and F.~Luca (eds.), \emph{Anatomy of Integers}, CRM Proceedings and Lecture Notes 46, American Mathematical Society, 2008. G.~Tenenbaum, \emph{Introduction to Analytic and Probabilistic Number Theory}, 3rd ed., AMS, 2015, Chapters~III.5 and~III.6.
\end{itemize}

